\section{Zusammenfassung und Ausblick}
"`Ziel des Praktikums ist es, die in einem Studium erworbenen theoretischen Kenntnisse im praktischen Umfeld anzuwenden und individuelle F�higkeiten im Umgang mit diesen zu erwerben. Der Praktikumsbericht enth�lt daher eine ausf�hrliche Reflexion der im Praktikum durchgef�hrten T�tigkeiten und setzt diese in Bezug zu den theoretischen Kenntnissen, welche im Studium erlangt wurden"' (\cite{Balz2008}, S. 158).

Der Schlussteil besteht aus einer kurzen Zusammenfassung und kann auch pers�nliche Bemerkungen, Bewertungen etc. enthalten. Zuerst wird somit das Vorher geschriebene knapp zusammengefasst. Anschlie�end werde die Ergebnisse reflektiert. Die Zusammenfassung setzt sich kritisch mit der Aufgabenstellung auseinander und beantwortet die Fragen:
\begin{itemize}
	\item Was sollte getan werden und was wurde tats�chlich erreicht?
	\item Wo gab es besondere Schwierigkeiten?
	\item Welche Teile der Aufgabenstellung konnten eventuell gar nicht bearbeitet werden und warum?
\end{itemize}

Der Ausblick kann Empfehlungen f�r eine weitere sinnvolle Fortsetzung des Projektes bzw. der Forschungsarbeit enthalten.
Solche Empfehlungen beruhen �blicherweise auf Erkenntnissen, die w�hrend der Bearbeitung des laufenden Projektes entstanden sind (vgl. \cite{Hering2009}, S. 44).\footnote{Es wird darauf hingewiesen, dass die wissenschaftliche Arbeit am Ende unbedingt von einer dritten Person korrekturgelesen werden sollte.
Dadurch k�nnen Fehlerarten wie falsche Wortwahl, Rechtschreibfehler. Zeichensetzungsfehler, Grammatikfehler, Stilfehler, inhaltliche Wiederholungen, Fehler im logischen Gedankenfluss, Layout-Fehler sowie sonstige Fehler aufgedeckt werden (vgl. \cite{Hering2009}, S. 174).
Zudem wird der Autor beim h�ufigen Lesen der eigenen Texte mit der Zeit "`betriebsblind"' gegen�ber den eigenen Formulierungen.
Auch aus diesem Grund sollte von einem Au�enstehenden �berpr�ft werden, ob der Bericht f�r Projektfremde verst�ndlich ist (vgl. \cite{Hering2009}, S. 130).}